
\subsubsection{6.1 Interpretation}

The results suggest that benchmark clustering reflects distinct task categories. Factual Recall benchmarks (TriviaQA, NQ, SQuAD) probe knowledge retrieval from parametric memory; Entity/Claim benchmarks (PopQA, HaluEval, FEVER) test long-tail entity knowledge and evidence-based verification. The approximately $7\times$ transfer gap between within-cluster and cross-cluster pairs indicates that distribution similarity is a relevant factor in transfer success.

\subsubsection{6.2 Practical Implications}

For practitioners deploying uncertainty-based hallucination detectors, these results suggest: (1) compute JS-divergence between source and target benchmark entropy distributions, (2) if distributions are similar (same cluster), threshold reuse may be feasible with limited degradation (~3\%), (3) if distributions differ (cross-cluster), recalibration is likely necessary (>20\% degradation observed here).

\subsubsection{6.3 Limitations}

\textbf{Model scope:} Only Llama-2-7B-Chat was evaluated. Larger models may exhibit different clustering structure and transfer behavior.

\textbf{Task scope:} Only short-form factual QA was tested. Summarization and long-form generation remain untested.

\textbf{Scale:} Proof-of-concept experiments used 100-1000 samples per benchmark. Full-scale validation with larger samples is needed.

\textbf{Experiment type:} H-E1 clustering used synthetic entropy distributions modeled on expected benchmark characteristics. H-M3 tested only one within-cluster pair (trivia\_qa $\leftrightarrow$ squad). H-M4 tested two cross-cluster pairs.

\textbf{Cross-cluster validation:} H-M4 cross-cluster degradation estimates were derived from a limited set of transfer pairs. The $7\times$ gap is observed in the tested pairs but requires validation across all cross-cluster combinations.

\textbf{Statistical power:} The Mann-Whitney comparison of within-cluster (n=2) versus cross-cluster (n=2) degradation values in H-M4 did not reach statistical significance (p=0.110) due to small sample sizes, though the effect size was large.

